% GENERATED FILE. Edit canonical lessons, then run npm run generate:books.
% canonical-source-hash: fnv1a64:768af7958201a1ea
% canonical-lessons: RU-C211-myshtsa, RU-C211-kost, RU-C211-puls, RU-C211-davleniye, RU-C211-sila

\chapter{Muscle, Bone, Pulse, Pressure, Strength}
\label{ch:ru-muscle-bone-pulse-pressure-strength}
\hlchaptermodality{211}

\begin{chapteropening}
\textbf{By the end of this chapter:} \emph{I can say a muscle, a bone, a pulse, pressure and strength.}

Complete the last lesson of chapter 211: I can say a muscle, a bone, a pulse, pressure and strength.
\end{chapteropening}

\section[\texorpdfstring{\ru{мышца} (mýshtsa)}{мышца (mýshtsa)}]{\ru{мышца} (mýshtsa) — a muscle}

\label{lesson:RU-C211-myshtsa}



\begin{quote}
Before the new one: say the Russian for a sticking plaster, then the Russian for a bandage.
\end{quote}

\subsection*{You'll want to know: \ru{мышца}}
\textbf{\ru{мышца}} (\emph{mýshtsa}) — \textquotedblleft{}a muscle\textquotedblright{}.

\begin{grammarlens}[title={What you've built}]
One more word for this chapter.
\end{grammarlens}

\subsection*{Your turn}
\begin{itemize}
\raggedright
  \item \emph{Say it:} \emph{mýshtsa}
  \item \emph{Say it:} \emph{mýshtsa}, once more
  \item \emph{Say it:} say \emph{mýshtsa}
  \item \emph{Recall:} say the Russian for a sticking plaster, then the Russian for a bandage, then say \emph{mýshtsa} again
\end{itemize}

\begin{tcolorbox}[breakable,colback=teal!4,colframe=teal!35!black,title={Before you move on}]
What does \ru{мышца} mean? (\textquotedblleft{}A muscle\textquotedblright{}.) Say it once more.
\end{tcolorbox}
\section[\texorpdfstring{\ru{кость} (kost)}{кость (kost)}]{\ru{кость} (kost) — a bone}

\label{lesson:RU-C211-kost}



\begin{quote}
Before the new one: say the Russian for a bandage, then the Russian for a muscle.
\end{quote}

\subsection*{You'll want to know: \ru{кость}}
\textbf{\ru{кость}} (\emph{kost}) — \textquotedblleft{}a bone\textquotedblright{}.

\begin{grammarlens}[title={What you've built}]
One more word for this chapter.
\end{grammarlens}

\subsection*{Your turn}
\begin{itemize}
\raggedright
  \item \emph{Say it:} \emph{kost}
  \item \emph{Say it:} \emph{kost}, once more
  \item \emph{Say it:} say \emph{kost}
  \item \emph{Recall:} say the Russian for a bandage, then the Russian for a muscle, then say \emph{kost} again
\end{itemize}

\begin{tcolorbox}[breakable,colback=teal!4,colframe=teal!35!black,title={Before you move on}]
What does \ru{кость} mean? (\textquotedblleft{}A bone\textquotedblright{}.) Say it once more.
\end{tcolorbox}
\section[\texorpdfstring{\ru{пульс} (puls)}{пульс (puls)}]{\ru{пульс} (puls) — a pulse}

\label{lesson:RU-C211-puls}



\begin{quote}
Before the new one: say the Russian for a muscle, then the Russian for a bone.
\end{quote}

\subsection*{You'll want to know: \ru{пульс}}
\textbf{\ru{пульс}} (\emph{puls}) — \textquotedblleft{}a pulse\textquotedblright{}.

\begin{grammarlens}[title={What you've built}]
One more word for this chapter.
\end{grammarlens}

\subsection*{Your turn}
\begin{itemize}
\raggedright
  \item \emph{Say it:} \emph{puls}
  \item \emph{Say it:} \emph{puls}, once more
  \item \emph{Say it:} say \emph{puls}
  \item \emph{Recall:} say the Russian for a muscle, then the Russian for a bone, then say \emph{puls} again
\end{itemize}

\begin{tcolorbox}[breakable,colback=teal!4,colframe=teal!35!black,title={Before you move on}]
What does \ru{пульс} mean? (\textquotedblleft{}A pulse\textquotedblright{}.) Say it once more.
\end{tcolorbox}
\section[\texorpdfstring{\ru{давление} (davléniye)}{давление (davléniye)}]{\ru{давление} (davléniye) — pressure; blood pressure}

\label{lesson:RU-C211-davleniye}



\begin{quote}
Before the new one: say the Russian for a bone, then the Russian for a pulse.
\end{quote}

\subsection*{You'll want to know: \ru{давление}}
\textbf{\ru{давление}} (\emph{davléniye}) — \textquotedblleft{}pressure; blood pressure\textquotedblright{}.

\begin{grammarlens}[title={What you've built}]
One more word for this chapter.
\end{grammarlens}

\subsection*{Your turn}
\begin{itemize}
\raggedright
  \item \emph{Say it:} \emph{davléniye}
  \item \emph{Say it:} \emph{davléniye}, once more
  \item \emph{Say it:} say \emph{davléniye}
  \item \emph{Recall:} say the Russian for a bone, then the Russian for a pulse, then say \emph{davléniye} again
\end{itemize}

\begin{tcolorbox}[breakable,colback=teal!4,colframe=teal!35!black,title={Before you move on}]
What does \ru{давление} mean? (\textquotedblleft{}Pressure; blood pressure\textquotedblright{}.) Say it once more.
\end{tcolorbox}
\section[\texorpdfstring{\ru{сила} (síla)}{сила (síla)}]{\ru{сила} (síla) — strength}

\label{lesson:RU-C211-sila}



\begin{quote}
Before the new one: say the Russian for a pulse, then the Russian for pressure; blood pressure.
\end{quote}

\subsection*{You'll want to know: \ru{сила}}
\textbf{\ru{сила}} (\emph{síla}) — \textquotedblleft{}strength\textquotedblright{}.

\begin{grammarlens}[title={What you've built}]
One more word for this chapter.
\end{grammarlens}

\subsection*{Your turn}
\begin{itemize}
\raggedright
  \item \emph{Say it:} \emph{síla}
  \item \emph{Say it:} \emph{síla}, once more
  \item \emph{Say it:} say \emph{síla}
  \item \emph{Recall:} say the Russian for a pulse, then the Russian for pressure; blood pressure, then say \emph{síla} again
\end{itemize}

\begin{tcolorbox}[breakable,colback=teal!4,colframe=teal!35!black,title={Before you move on}]
What does \ru{сила} mean? (\textquotedblleft{}Strength\textquotedblright{}.) Say it once more.
\end{tcolorbox}
